\tikzset{
  host/.style={
    scale=\nodescale,
    inner sep=1pt,
    draw,
    rounded rectangle,
  },
  edge/.style={
    scale=\nodescale,
    inner sep=1pt,
    draw,
    rounded rectangle,
    fill=blue!20,
  },
  core/.style={
    scale=\nodescale,
    inner sep=1pt,
    draw,
    rounded rectangle,
    fill=green!30,
  },
  port/.style = {
    % font=\tiny,
    scale=.5,
    inner sep=1pt,
    draw,
    signal,
    fill=white,
  }
}
\newcounter{previtem}

\section{Implementing and Testing a Probing-Based Path Verification (PBV)}

PBV consists of sending special packets (probes) with path-verification capabilities. The path verification system used is described \autoref{sec:multisigmodel}. Some assumptions are made:
\begin{inlinelist}
    \item Each node is assumed to be secure, that is, no node will alter the packet in any way that is not expected. This is a common assumption in SDN networks, where a trusted party is the only entity that can alter the network state;
    \item  Every link is assumed to be perfect, that is, no packet loss, no packet corruption, and no packet duplication;
    \item  Protocol boundary is IPv4. This means that \polka is only used inside this network, and only IPv4 is used outside;
    \item All paths are assumed to be valid and all information correct unless stated otherwise;
    \item Hash collisions will not happen.
\end{inlinelist}
Furthermore, this implementation is a proof-of-concept\mycomment{está certo chamar de proof-of-concept?}, and there is little concerns about performance and protocol implementation: only IPv4 is supported. The following subsections detail the setup and implementation.

\subsection{Setup}

\pIV is a language used to program the data plane of network devices\cite{p4}. Behavioral Model 2 (\bmv) is a software switch that supports the \pIV language \cite{bmv2}. Mininet is a network emulation tool that allows the creation of virtual networks with a controlled environment\cite{mininet}. Mininet-wifi\footnote{\url{github.com/intrig-unicamp/mininet-wifi}} is a Mininet fork used in this work due to its compatibility with \bmv. Despite the name, no Wi-Fi emulating feature was used.

% Explain why using Scapy
Scapy is a Python library that allows the creation and manipulation of packets\cite{scapy}. It is used to parse packets in the experiments and automate tests. The setup is shown in \autoref{setup}. Controller is responsible for the network state and reference signature calculation, and the Tester triggering Mininet hosts to send packets, setup Scapy for parsing them and then asserting their signatures to Controller-calculated references.



\subsection{Implementation}

We define the header fields \timestamp and \lhash to indicate the uniquely-generated number and the current digest, respectively. Due to header incompatibility, probe packets operate using a different version number in the \texttt{version} field, so it is interoperable with \polka by switching between probes and regular packets using the \texttt{version} header. \polka packets use version \texttt{0x01}, and probe packets use version \texttt{0xF1}. \autoref{topology:base} shows the used topology used in the experiments where $s_i$ is a switch, $e_i$ is an edge node, and $h_i$ is a host. The topology is a linear chain with 10 switches and 11 hosts. The links are bidirectional, but only one direction is shown for simplicity. %The ports are numbered from $1$ to $3$, and the ports used are shown in the figure. The topology is used to test the path verification mechanism.

\begin{figure}[ht]
    \centering    
    \begin{minipage}{0.49\linewidth}
    \centering    
    \begin{tikzpicture}[
        outer sep=auto,
        every node/.style={draw, rectangle, label distance=-4pt},
        switch/.style={chamfered rectangle, fill=green!20, inner sep=1pt},
        % every fit/.style={rectangle,draw}
    ]
        \node (controller) {Controller};
        \node[right=12pt of controller] (tester) {Tester};
        \draw[->] (controller) to (tester);
        \node[label=above:Python script, dotted,
          fit={(controller) (tester)}] (py) {};

        \node[switch, below=32pt] at (py) (s1) {cores};
        \node[switch, left=12pt] at (s1.west) (ingress) {ingress};
        \node[switch, right=12pt] at (s1.east)      (egress) {egress};
        \draw (ingress) to (s1);
        \draw (s1) to (egress);
        \node[label=below:Mininet, dotted,
            fit={(ingress) (s1) (egress)}] {};

        \draw[->] (ingress) -- node[fill=white, inner sep=1pt] {\footnotesize Scapy} (controller);
        \draw[->] (egress) -- node[fill=white, inner sep=1pt] {\footnotesize Scapy} (tester) ;
    \end{tikzpicture}
    \caption{Experimentation setup}
    \label{setup}
    \end{minipage}
    \begin{minipage}{0.49\linewidth}
    \centering    
    \newcommand{\nodescale}{0.8}
    \newcommand{\nswitches}{10}
    % with ports
    % \newcommand{\hspacing}{0.95}
    % \newcommand{\vspacing}{0.95}
    % \newcommand{\nodescale}{0.8}
    % \newcommand{\nswitches}{10}
    % \begin{tikzpicture}[outer sep=auto]
    %     \foreach \x in {1, ..., \nswitches} {
    %         \node[host] (H\x) at (\x*\hspacing, 0*\vspacing) {$h\x$};
    %         \node[edge] (E\x) at (\x*\hspacing, 0.95*\vspacing) {$e\x$};
    %         \node[core] (S\x) at (\x*\hspacing, 2*\vspacing) {$s\x$};

    %         \draw (H\x) to
    %             node[pos=0.1, port, signal to=north]{1}
    %             node[pos=0.9, port, signal to=south]{1}
    %             (E\x);
    %         \draw (E\x) to
    %             node[pos=0.1, port, signal to=north]{2}
    %             node[pos=0.9, port, signal to=south]{1}
    %             (S\x);
    %     };
    %     \node[host] (H11) at (0*\hspacing, 0*\vspacing) {$h11$};
    %     \draw (H11) to
    %         node[pos=0.1, port, signal to=north]{1}
    %         node[pos=0.9, port, signal to=south]{3}
    %         (E1);
        
    %     \draw (S1) to
    %         node[pos=-0.1, port, signal to=right]{2}
    %         node[pos=1.1, port, signal to=left]{2}
    %         (S2);
    %     \foreach \x in {2, ...,  \number\numexpr\nswitches-1\relax } {
    %         \draw (S\x) to
    %             node[pos=-0.1, port, signal to=right]{3}
    %             node[pos=1.1, port, signal to=left]{2}
    %             (S\number\numexpr\x+ 1\relax);
    %     };
            
    % \end{tikzpicture}

    % Without ports

    \newcommand{\newtopoitem}[2]{%
        \node[host, right=2pt] (H#1) at (H#2.east)  {$h_{#1}$};
        \node[edge, above=4pt] (E#1) at (H#1.north) {$e_{#1}$};
        \node[core, above=4pt] (S#1) at (E#1.north) {$s_{#1}$};
        \draw (H#1) -- (E#1) -- (S#1);
    }
    \setcounter{previtem}{11}
    \begin{tikzpicture}[outer sep=auto]
        \node[host] (H11) {$h_11$};
        \foreach \x in {1, ..., \nswitches} {
            \newtopoitem{\x}{\arabic{previtem}}
            \ifnum \theprevitem < 11
                \draw (S\x) to (S\theprevitem);
            \fi
            \setcounter{previtem}{\x}
        };
        \draw (H11) to (E1);
            
    \end{tikzpicture}
    \caption{Baseline topology used in experimentation}
    \label{topology:base}
    \end{minipage}
% \end{figure}
% Joins both figures into one
% \begin{figure}[ht]
%     \centering
    \tikzset{
        outer sep=auto,
        every node/.style={draw, chamfered rectangle, label distance=-4pt, inner sep=2pt},
    }
    \begin{minipage}{0.49\linewidth}
    \centering    
    \begin{tikzpicture}
        \node (parse) {Parse};
            
        \node[above right=24pt] (encap) at (parse.east) {Encapsulation};
        \node[right=12pt] (core) at (encap.east) {SDN core};
        \draw[->, very thick] (parse.east) -- (encap)
            node[draw=none,pos=0.7,xshift=2pt,anchor=mid west] {\footnotesize if \texttt{0x0800} (ethernet)};
        \draw[->, very thick] (encap) -- (core);
        
        \node[below right=24pt] (decap) at (parse.east) {Decapsulation};
        \node[right=12pt] (out) at (decap.east) {external net};
        \draw[->, very thick] (parse.east) -- (decap)
            node[draw=none,pos=0.7,xshift=2pt, anchor=mid west] {\footnotesize if \texttt{0x1234} (polka)};
        \draw[->, very thick] (decap) -- (out);
    \end{tikzpicture}
    \caption{Edge ingress/egress steps}
    \label{fig:pipeline-edge}
    \end{minipage}
    \begin{minipage}{0.49\linewidth}
    \centering    
    \begin{tikzpicture}
        \node (parse) {Parse};
        \node[right=12pt] (nhop) at (parse.east) {nHop};
        \node[right=12pt] (sign) at (nhop.east) {Sign};
        \node[right=12pt] (deparse) at (sign.east) {Deparse};

        \draw[->, very thick] (parse) -- (nhop);
        \draw[->, very thick] (nhop) -- (sign);
        \draw[->, very thick] 
            (nhop.east) ..controls ++(0,-1) and ++(0,-0.5) .. (deparse.195) node[midway,below=-4pt,draw=none] {\footnotesize if not probe};
        \draw[->, very thick] (sign) -- (deparse);
            
    \end{tikzpicture}
    \caption{Data plane steps \mycomment{Muito simples?}}
    \label{fig:pipeline-core}
    \end{minipage}
\end{figure}

On edge nodes, if an IPv4 protocol \texttt{ethertype} field is detected (\texttt{0x0800}), it must be a packet from outside the network, so the edge is doing an ingress and \polka headers need to be added. Let this process be called \textit{encapsulation}; If a \polka protocol \texttt{ethertype} field is detected (\texttt{0x1234}), it must be a packet from the network core, so the edge is doing an egress and the original IPv4 packet must be unwrapped. Let this process be called \textit{decapsulation}.

\polka headers consists of the route polynomial (\routeid), along with \texttt{version}, \texttt{ttl} and \texttt{proto} (stores the original \texttt{ethertype}). \routeid and hop calculation is out of the scope of this paper.
During encapsulation, the packet can be made into a probe packet by setting the \texttt{version} header field to \texttt{0xF1} and further defining the 32-bit fields \timestamp to any unique value and \lhash to \timestamp.

On the core network, the hash function used is \textit{SipHash-2-4-32} (\siphash), which after calculating the next hop (\nhop) and if the packet is a probe packet \lhash is computed as
$$
\lhash \leftarrow \textit{SipHash2-4-32}(\nodeid || \nhop || \timestamp || 0000000_2, \lhash)
$$
Meaning that a 64-bitstring made of concatenating \nodeid, \nhop, \timestamp and 7 bit wide zero is used as key for the \siphash function hashing \lhash. Inputting \nhop into the hash function protects us from some scenarios. \mycomment{explicar mais}
